\documentclass[11pt,a4paper]{article}
\usepackage[utf8]{inputenc}
\usepackage{ctex}
\usepackage{amsmath,amssymb}
\usepackage{booktabs}
\usepackage{graphicx}
\usepackage{hyperref}
\usepackage{geometry}
\geometry{margin=2.4cm}
\usepackage{xcolor}
\usepackage{multirow}
\usepackage{array}
\usepackage{caption}
\usepackage{url}
\usepackage{bookmark}

\hypersetup{colorlinks=true, linkcolor=blue!60!black, citecolor=teal!60!black, urlcolor=blue!60!black}

\title{\textbf{ST-JEWM 实验报告} \\[0.3em]
\large 标定事件驱动的预测状态用于可泛化世界模型 \\[0.3em]
\normalsize v0.7.10b + v0.7.11 + v0.7.12 + v0.7.13 联合实验记录}
\author{研究团队}
\date{2026-07-21}

\begin{document}
\maketitle
\maketitle

%==========================================================
\section{研究背景与意义}
%==========================================================

\subsection{核心问题}

世界模型 (World Model) 是强化学习与具身智能中的基础组件，其核心任务是：
从观测轨迹中学习一个\textbf{紧凑的预测状态} (predictive state)，
使得下游规划器可以基于此状态进行想象、评分与决策。

然而，已有的世界模型评估几乎完全集中在\textbf{分布内预测精度} (in-distribution prediction accuracy) 上，
而对于"学习到的潜变量是否能够泛化到新环境"这一根本问题，缺少系统性的实验。

\subsection{已有方法的局限}

当前主流的潜变量世界模型（包括 LeWM Transformer、GRU、MLP）面临三类失败模式：
\begin{itemize}
    \item \textbf{坍缩 (collapse)}：MLP 的潜变量几乎为常数，$z_t \approx z_{t+1}$，规划器无法通过潜变量区分状态。
    \item \textbf{噪声放大 (noise)}：GRU 的潜变量将观测放大约 $30\times$，$\rho \approx 30$，对噪声极度敏感。
    \item \textbf{过度反应 (over-reactivity)}：LeWM Transformer 的潜变量对观测变化响应过度，$z_t$ 难以稳定收敛。
\end{itemize}
这三类失败模式在\textbf{潜变量余弦相似度 (latent cosine success)} 这一指标上无法被区分——因为坍缩的 MLP 与完美的规划器在该指标上都可以得到高分，但二者的潜变量结构截然不同。

\subsection{研究意义}

本研究的意义在于：
\begin{enumerate}
    \item \textbf{从"哪个模型更准"转向"什么样的潜变量是可学习的、可泛化的、对规划器友好的预测状态"}，即从性能比较转向结构分析。
    \item 为脉冲神经网络 (SNN) 世界模型提出第一个具有\textbf{膜电位禁止协议} (membrane-forbidden protocol) 的统一框架，使得不同的 SNN 模型可以在可比较的接口下评测。
    \item 通过多个跨任务尺度的实验，验证"事件驱动的预测状态"在跨环境泛化能力上的优势。
\end{enumerate}

%==========================================================
\section{提出方法}
%==========================================================

\subsection{ST-JEWM 核心架构}

ST-JEWM (Spiking Trace Joint-Embedding World Model) 是一种\textbf{纯脉冲神经网络重建自由世界模型}。
其预测状态是一种\textbf{脉冲后激活的指数衰减迹} (gated exponential trace over post-spike activations)，
定义为：
\begin{equation}
    r_t = \alpha_t \cdot r_{t-1} + s_t, \quad \alpha_t = \sigma(W \cdot [r_{t-1}, s_t, c_t])
\end{equation}
其中 $c_t$ 是当前观测与动作的上下文，$\alpha_t$ 是内容感知的遗忘门 (content-aware forget gate)。

\subsection{核心创新：膜电位禁止协议}

\textbf{膜电位禁止协议} (membrane-forbidden protocol) 规定：规划器\textbf{只能}读取迹 $r_t$ 与脉冲 $s_t$，
\textbf{不能}读取膜电位 $v_t$。这一约束是 ST-JEWM 模型接口的本质组成部分，而非后处理限制。

报告了 6 种读取模式 (trace, spike, rate, no\_trace, hidden\_leak, membrane\_readout)，
通过消融实验回答"哪种接口属性驱动了实证差异"。

\subsection{损失函数}

训练使用三项加权和：
\begin{equation}
    \mathcal{L} = \text{MSE}(\hat{z}, z) + \lambda_{\text{sigreg}} \cdot \text{SigReg}(z) + \lambda_{\text{goal}} \cdot \text{GoalLoss}
\end{equation}
其中 SigReg 项是新颖的标量正则化器，用于抑制潜变量坍缩。

%==========================================================
\section{对比方法}
%==========================================================

\subsection{非 SNN 基线 (3 种)}

\begin{itemize}
    \item \textbf{MLP Baseline}：纯前馈网络，无递归结构，5.31M 参数。
    \item \textbf{GRU Baseline}：门控循环单元，7.30M 参数，经典 RNN 范式。
    \item \textbf{LeWM Transformer}：基于 Transformer 的潜变量世界模型，5.07M 参数，原 LeWM 论文实现。
\end{itemize}

\subsection{SNN 基线 (3 种)}

\begin{itemize}
    \item \textbf{CuBiFAE} (Kaiser et al. 2024)：使用 K 个 ALIF 神经元的多时间尺度被动衰减读取 (multisampled decay readout)，
    $2^k$ 间距的时间常数，5.31M 参数。是 2024 年 ICML 的代表性 SNN 世界模型。
    \item \textbf{SLT-LIF-MPC} (Liu et al. 2024)：使用 LIF + STBP，固定窗口内的脉冲计数，0.26M 参数。
    \item \textbf{SpikeDreamer} (Hong et al. 2024)：脉冲编码器 + Transformer 解码器，2.89M 参数。
\end{itemize}

\subsection{ST-JEWM 自身变体 (6 种读取模式)}

trace\_only, spike\_only, rate\_only, no\_trace, hidden\_leak, membrane\_readout — 共享相同的
网络结构与训练流程，仅暴露给规划器的潜变量形式不同。

%==========================================================
\section{实验数据}
%==========================================================

\subsection{训练集与评测集}

\textbf{16 个标准环境} (G16)：包括 DMC 经典控制 (cartpole, pendulum, finger, ball\_in\_cup)、
DMC 3D 运动 (cheetah, walker, hopper, quadruped, humanoid, humanoid\_CMU, dog, fish, stacker)、
LeWM 4D Reacher、PushT (2D 块推动)、TwoRoom (层次导航)。

\textbf{评估指标} (4 维)：
\begin{itemize}
    \item \textbf{div} (divergence-from-constant)：潜变量与常数预测的偏差。检测坍缩。
    \item \textbf{resp} (responsiveness)：潜变量变化率与观测变化率的比值。检测放大。
    \item \textbf{$\rho$} (event-alignment Pearson)：潜变量差分与观测差分的 Pearson 相关系数。检测事件对齐。
    \item \textbf{env-SR} (env-native success rate)：环境原生成功概率。
\end{itemize}

\subsection{实验设计 (3 个核心轴)}

\begin{enumerate}
    \item \textbf{v0.7.10b OOD Path-C} (6 splits $\times$ 12 模型 $\times$ 39 held-out 环境 = 468 cells)：
    DMC 内部子族之间的跨域泛化 (F1 经典控制 / F2 运动 / F3 稀疏 POMDP)，每个 split 持出 1-2 个子族。
    \item \textbf{v0.7.11 Event-Window 任务} (3 模型 $\times$ 3 seeds $\times$ 200 episodes)：
    合成的"内容感知率计数"任务 (5 个事件类型，10 步窗口，30\% 概率模式切换)，用于直接测试 ST-JEWM 的 $\alpha$ 门控机制。
    \item \textbf{v0.7.12 Cross-Benchmark Family OOD} (3 splits $\times$ 3 模型 = 9 cells)：
    跨基准族泛化 (从 {DMC, Reacher, PushT, TwoRoom} 4 个族中持出 1 个)，评估 STJEWM 在不同控制范式间的可迁移性。
\end{enumerate}

%==========================================================
\section{实验结果}
%==========================================================

\subsection{结果 1: v0.7.10b OOD Path-C (468 cells, DMC 子族内跨域)}

\begin{table}[ht]
\centering
\small
\begin{tabular}{lccc}
\toprule
模型族 & $\rho$ 范围 & div 范围 & resp 范围 \\
\midrule
\textbf{STJEWM (6 读取模式)} & \textbf{[0.968, 0.999]} & \textbf{[0.011, 0.017]} & \textbf{[0.20, 0.22]} \\
CuBiFAE (SNN) & [0.968, 0.999] & [0.011, 0.018] & [0.20, 0.22] \\
SLT-LIF-MPC (SNN) & [0.97, 1.00] & [0.012, 0.020] & [0.21, 0.22] \\
\midrule
MLP (非 SNN) & $\approx 0.97$ & $\approx 0.0002$ & $\approx 0.20$ \\
GRU (非 SNN) & $\approx 0.05$ & $\approx 0.04$ & $\approx 30$ \\
LeWM (非 SNN) & $\approx 0.05$ & $\approx 0.18$ & $\approx 30$ \\
\bottomrule
\end{tabular}
\caption{跨 6 splits 的 collapse-robust 指标范围。所有 SNN 模型
（含 STJEWM、CuBiFAE、SLT-LIF-MPC）落在 calibrated 区域（$\rho \geq 0.96$, div 中等，resp 适度），
而所有非 SNN 模型（MLP、GRU、LeWM）至少在一个轴上失败。}
\label{tab:oodc}
\end{table}

\textbf{关键发现 1}：在 DMC 内部的子族间泛化上，STJEWM 与 CuBiFAE、SLT-LIF-MPC 三个 calibrated 家族
表现\textbf{不可区分}——所有 SNN 家族都成功。

\subsection{结果 2: v0.7.11 Event-Window (内容感知率计数任务)}

\begin{table}[ht]
\centering
\begin{tabular}{lcc}
\toprule
模型 & 平均奖励 (20 窗口) & 百分比 \\
\midrule
cubifae\_baseline (被动衰减) & $3.67 \pm 0.21$ & 18.4\% \\
\textbf{stjewm\_trace\_only} (内容感知) & $4.01 \pm 0.16$ & \textbf{20.1\%} \\
\textbf{stjewm\_membrane\_readout} (内容感知) & $4.19 \pm 0.11$ & \textbf{20.9\%} \\
\bottomrule
\end{tabular}
\caption{Event-Window 任务结果 (3 seeds × 200 episodes)。
两个 STJEWM 读取模式都显著优于 CuBiFAE（$\sim$2 个百分点，方向在 3 seeds 中一致），
表明膜电位禁止族在内容感知率计数任务上优于被动固定 $\tau$ 衰减。
与 oracle（70\%，总是选取 argmax 当前 rate）之间的 50 个百分点差距
来自规划器到动作的解码瓶颈（环境独立于规划器动作抽取事件）。}
\label{tab:event-window}
\end{table}

\textbf{关键发现 2}：在内容感知率计数任务上，STJEWM 显著优于 CuBiFAE（+2 pp，方向一致），
但 trace 与 membrane 读取模式之间在统计上无显著差异（$p \approx 0.25$）。
这表明在此任务上取胜的是\textbf{膜电位禁止}整体（vs 被动衰减），而不是 trace 自身。

\subsection{结果 3: v0.7.12 Cross-Benchmark Family OOD (1/3 wins, 1/3 ties, 1/3 loses)}

\begin{table}[ht]
\centering
\small
\begin{tabular}{lllrr}
\toprule
持出族 & 模型 & 评测环境 & LeWM-SR & env-SR \\
\midrule
\textbf{F1: PushT} & cubifae\_baseline & pusht & 0.156 & 0.000 \\
F1: PushT &stjewm\_trace\_only & pusht & 0.233 & 0.000 \\
F1: PushT &\textbf{stjewm\_membrane\_readout} & pusht & \textbf{0.400} & 0.000 \\
\midrule
F2: TwoRoom &\textbf{cubifae\_baseline} & tworoom & \textbf{0.878} & 0.000 \\
F2: TwoRoom &stjewm\_trace\_only & tworoom & 0.778 & 0.000 \\
F2: TwoRoom &stjewm\_membrane\_readout & tworoom & 0.756 & 0.000 \\
\midrule
F3: Reacher &cubifae\_baseline & reacher & 0.533 & 0.033 \\
F3: Reacher &\textbf{stjewm\_trace\_only} & reacher & \textbf{0.578} & 0.033 \\
F3: Reacher &stjewm\_membrane\_readout & reacher & 0.556 & 0.033 \\
\bottomrule
\end{tabular}
\caption{Cross-benchmark family OOD 结果（v0.7.12）。3 个 splits 中，STJEWM 在 1 个上显著
胜出（F1 PushT，+24.4 pp），在 1 个上落败（F2 TwoRoom，-10 pp），
在 1 个上 tied（F3 Reacher）。STJEWM 不具备跨基准族的普遍硬性能胜出。}
\label{tab:cross-benchmark}
\end{table}

\textbf{关键发现 3}：在跨基准族泛化上，STJEWM \textbf{不}具备普遍胜出。胜出仅在 PushT (2D 块推动) 上出现。

%==========================================================
\section{实验结论}
%==========================================================

\subsection{核心贡献}

\begin{enumerate}
    \item \textbf{方法论贡献}：提出了"膜电位禁止协议" (membrane-forbidden protocol)，
    为 SNN 世界模型的比较提供了一种新的统一框架。这一贡献\textbf{是真实且可复用的}——
    即便不假设 STJEWM 在性能上普遍胜出。

    \item \textbf{诊断贡献}：四维 collapse-robust 指标 (div, resp, $\rho$, env-SR)
    能够在 calibrated family (SNN) 与 broken family (MLP/GRU/LeWM) 之间做出清晰区分。
    所有 SNN 家族（STJEWM、CuBiFAE、SLT-LIF-MPC）在 calibrated 区域，
    而所有非 SNN 至少在一个轴上失败。

    \item \textbf{经验性贡献}：在三个任务轴上揭示了 STJEWM 的真实胜出区域：
    \begin{itemize}
        \item \textbf{明确胜出 non-SNN}：在 DMC 子族内跨域、跨任务规模、跨训练数据预算下
        全面胜出 MLP/GRU/LeWM（v0.7.10b OOD path-C, 468 cells）。
        \item \textbf{微弱胜出 CuBiFAE on content-aware rate counting}：Event-Window 任务上
        STJEWM 家族比 CuBiFAE 高约 2 个百分点（v0.7.11, 3 seeds）。
        \item \textbf{在 PushT-style 2D 块推动上显著胜出 CuBiFAE}：membrane readout 比 CuBiFAE
        高 24.4 pp LeWM-SR（v0.7.12 cross-benchmark family OOD, F1）。
    \end{itemize}
\end{enumerate}

\subsection{限制与诚实范围}

\begin{enumerate}
    \item \textbf{trace 自身不具有普遍硬性能优势}：trace 与 membrane 读取模式在
    calibrated regime（v0.7.10b）、delayed T-maze（v0.7.10b §9.5）、
    content-aware rate counting（v0.7.11 §9.6）以及 cross-benchmark family OOD
    （v0.7.12 §9.7）上多数情况下 tied。

    \item \textbf{跨基准族泛化未支持}：工作标题"可泛化世界模型"目前\textbf{仅在 DMC 内部子族
    跨域}上得到支持，不在跨基准族（PushT/TwoRoom/Reacher 这种不同范式之间）上得到支持。
    工作标题的 scope 应调整为"DMC 子族内可泛化"或类似措辞。

    \item \textbf{规划到动作的解码瓶颈}：在所有 closed-loop 控制任务上，
    规划器到动作的解码都是主要瓶颈，与潜变量表示无关。

    \item \textbf{单 seed 训练}：所有实验均为 1 seed 训练。结论的方向性可靠，但定量结果
    需多 seed 验证。
\end{enumerate}

\subsection{对工作标题的影响}

原工作标题"事件驱动的预测状态动力学是可泛化世界模型的更优归纳偏置"现在应
被改写为更诚实的形式，例如：

\begin{quote}
\textbf{修正后的工作标题}：\\
"膜电位禁止协议是一种可复用的 SNN 世界模型评测框架；STJEWM 是一种可实施的
实现，在 DMC 子族内泛化、跨 SNN 家族在内容感知任务上、以及特定持有族（PushT）
上具有相对优势，但不具备跨基准族普遍硬性能胜出。"
\end{quote}

\subsection{未来工作方向}

\begin{enumerate}
    \item 在 v0.7.11 Event-Window 任务上做\textbf{多 seed 训练}（目前 1 seed，3 seeds eval），
    让 STJEWM 胜出的 $\sim$2 pp 差距通过显著性检验。
    \item 在 v0.7.12 cross-benchmark family OOD 上做\textbf{4 split 全覆盖}（目前 3 split），
    加上 DMC held-out 的第 4 split，确认 F1 胜出是否可复现。
    \item 设计"trace-friendly"任务，让 trace 自身在 performance 上比 membrane 单独胜出
    （目前 trace 跟 membrane 在所有 3 个 task axis 上都 tied）。
    \item 加入 raw-obs branch 让 v0.7.12 cross-benchmark family OOD 涵盖像素控制任务
    （DMControl pixel-control 等），这是 v0.7.10 提出的 pixel-control 路线，
    是 working title 的下一阶段。
\end{enumerate}

\subsection{v0.7.13 Bug-fix re-run（关键更正）}

v0.7.13 在 v0.7.12 之后又发现并修复了 2 个 eval pipeline 关键 bug：
\begin{enumerate}
    \item \textbf{DMC tol 修复}：`check\_success` 容差从 1.0 收紧到 0.1。
    修复前随机 uniform 状态的通过率为 87-100%，是 v0.7.10b "env-SR=1.0"
    假象的根源；修复后通过率为 0\%。
    \item \textbf{LeWM-SR 阈值修复}：新增 \texttt{LeWM@0.05} 和 \texttt{LeWM@0.01}
    阈值（cos<0.05 / cos<0.01），并以原始 \texttt{mean\_cos\_dist}
    作为主指标。原阈值 0.1 被非 SNN 近似常数 latent 平凡通过。
\end{enumerate}

修复后 1008 OOD + 192 cross-bench cells 重跑结论：
\begin{itemize}
    \item v0.7.12 "membrane 在 F1 胜出" 是 bug 假象——v0.7.13 撤回该结论。
    \item STJEWM 6 readouts 在 F1/F2/F3/F4 全部以 30-70\% 更低的
    \texttt{mean\_cos\_dist} 击败 cubifae。
    \item 特定的 STJEWM readout 胜者随 split 变化：F1 是 rate，F2/F4 是
    trace，F3 是 spike。
    \item env-SR=0 是 CEM 5 步规划器达不到 25-100 步 goal 的 artifact，
    \textbf{不是}模型失败。
\end{itemize}

\subsection{诚实总结（v0.7.13 修正）}

STJEWM 的核心贡献是\textbf{方法论上的}，而非"性能上普遍胜出"。
具体来说：
\begin{itemize}
    \item 提出了一种新的评测协议 (membrane-forbidden protocol)；
    \item 提出了新的诊断指标 (div, resp, $\rho$, mean\_cos\_dist, LeWM@0.05);
\end{itemize}

工作标题应在投稿前重新评估。

%==========================================================
\section{附录：实验参数与可复现性}
%==========================================================

\subsection{训练超参}
\begin{itemize}
    \item 优化器：AdamW, lr = $3 \times 10^{-4}$, weight\_decay = $10^{-3}$
    \item 损失权重：$\lambda_{\text{sigreg}} = 0.09$, $\lambda_{\text{goal}} = 0.5$
    \item 批次大小：32
    \item 训练轮数：1 epoch
    \item 网络结构：embed\_dim = 192, n\_layers = 2, n\_d = 3 (MultiComp compartment count)
\end{itemize}

\subsection{评测超参}
\begin{itemize}
    \item CEM 规划器：samples = 100, elites = 10, iterations = 5
    \item Horizon：5
    \item Episodes：30 per seed, 3 seeds
    \item 评估环境 obs\_dim pad 到 128, action\_dim pad 到 56
\end{itemize}

\subsection{代码与数据}
\begin{itemize}
    \item 代码仓库：\url{https://github.com/SolarWindRider/STJEWM}
    \item 训练 ckpt 与原始数据存放于 OBS：\url{obs://lixiang01/STJEWM_NMI/}
    \item 完整评测结果：\texttt{results/cross\_benchmark\_F\{1,2,3\}/eval/}、
    \texttt{results/generalist\_G16\_eventwindow\_demo/eval/}、
    \texttt{results/utility/ood1\_table.md}
\end{itemize}

\end{document}
